% Related work is defined here and placed where each connection becomes
% relevant in the five-section narrative.

\newcommand{\RelaxationTheoryPriorWork}{%
Two established ideas guide this interpretation. Fading memory asks whether
old inputs and the initial state gradually lose their influence
~\cite{JaegerESP2001,BoydChua1985,GrigoryevaOrtega2018}. Information-processing
capacity asks which functions of recent input history remain available to a
readout~\cite{Dambre2012,MartinezPenaIPC2023}. Quantum versions of these ideas
have been developed for finite, input-driven, and subsystem reservoirs
~\cite{MartinezPenaOrtega2023,MartinezPenaOrtega2025,KobayashiESP2024}.
Relaxation spectra provide a complementary fixed-input view of how an open
quantum system forgets~\cite{Minganti2018,Xia2026}. Because switched inputs mix
those responses, we use the spectral calculations as diagnostics rather than
as a complete mechanism.
}

\newcommand{\FiniteSamplePriorWork}{%
Finite-sample training, resolvable expressive capacity, and statistical-noise
robustness have been studied in QRC and physical learning systems
~\cite{Ahmed2024,HuSampling2023}. Our measurement study asks whether
the observed ordering survives when expectation values must be estimated from
the same nominal number of state preparations. It does not model a
platform-specific implementation cost.
}
